%% Full length research paper template
%% Created by Simon Hengchen and Nilo Pedrazzini for the Journal of Open Humanities Data (https://openhumanitiesdata.metajnl.com)

\documentclass[11pt]{article} 
\usepackage[english]{babel}
\usepackage[letterpaper]{geometry}
\usepackage{johd}
\usepackage{graphicx} 
\usepackage{xcolor}
\usepackage{amsmath, mathtools, mathrsfs, mathdots}
\usepackage{bm, amssymb, bbm}

\title{Marital Insurance and Job Search: Evidence from Property Regime Reform}

\author{Sofia Valdivia \& Heleen Ren}

\date{Feb. 2026}

\begin{document}

\maketitle

Individuals who marry gain the benefits of being in a two-person household—specialization within the household, dual contributions to public goods, and insurance in the face of future uncertainty. A central insight of the
household search literature is that these benefits lead married couples to behave differently from singles in the labor market. \cite{GULER2012352} develop a framework in which married households face joint location and employment decisions that generate both new opportunities and new frictions. In their model, reservation wages for employed-searcher couples can exceed those of singles because the spouse’s income provides a form of insurance. Building on this foundation, \cite{mankart2017household} show that household members adjust their labor supply to insure against unemployment risks, \cite{pilossoph_wee} demonstrate that income pooling raises reservation wages and accelerates job ladder climbing, and \cite{FERNANDEZBLANCO2022105477} establishes that spousal income simultaneously provides consumption insurance while increasing labor market risk-taking.

While this literature has established that insurance within marriage drives differences in job search behavior between singles and couples, it treats the strength of marital insurance as fixed. In practice, however, the extent of insurance within marriage varies considerably across couples depending on the legal and institutional design of the marital contract. Matrimonial property regimes are an important component of the marital contract. It determines how assets are owned during marriage and divided upon divorce. Two broad regimes exist: separation of property, under which each spouse retains ownership of their own assets, and community property, where assets acquired during marriage are jointly owned and divided equally. Within the latter, \textit{universal} community property extends joint ownership to all premarital assets and inheritances, whereas \textit{limited} community property includes only assets acquired after marriage. These rules govern how much risk is truly pooled inside the household and, importantly, what happens in the event of divorce.

% how sensitive are search decisions to the strength of the marital contract?
This paper tests whether the insurance provided by marriage shapes couples’ job search behavior. Specifically, we ask whether weakening marital insurance, through a weaker property regime or more individualized asset ownership, lowers reservation wages and reduces labor market risk-taking within couples. We build on the framework of \cite{GULER2012352}, in which spouses make joint search decisions and spousal income provides insurance, but depart from their approach by allowing the strength of that insurance to vary. Rather than comparing singles to couples, we compare couples facing different degrees of asset pooling and divorce protection. When the marriage contract weakens, because fewer assets are jointly owned or because divorce reallocates wealth differently, effective intra-household insurance declines. Even while married, the anticipation of divorce affects savings, labor supply, and job search decisions, bringing behavior closer to that of singles.


%This paper directly tests the insurance mechanism underlying differential job search behavior between singles and couples by exploiting variation in the strength of marital insurance. We ask whether variation in intra-household insurance changes equilibrium job search behavior within couples. In particular, does a reduction in marital insurance, induced by a weaker property regime or more individualized asset ownership, lower reservation wages and reduce labor market risk-taking? We use the \cite{GULER2012352} framework, where spouses make joint job search decisions based on the insurance provided by marriage, but we compare couples facing different degrees of marital insurance rather than comparing singles to couples. When the strength of the marriage contract declines, because fewer assets are jointly owned or because divorce reallocates assets differently, the amount of effective insurance inside marriage falls. Even while married, anticipation of potential divorce affects savings, labor supply, and job search behavior.

We exploit the 2018 Dutch reform that changed the default property regime from universal
to limited community property, thereby exogenously reducing the degree of financial insurance between
spouses. We combine this reform with variation in wealth shocks such as inheritances, gifts, and premarital asset appreciation. The reform re-defined which wealth shocks are shared and which remain individual. The data we use comes from the Dutch Household Survey (DHS), which records detailed information on household finances, marital status, and labor outcomes.

%Single individuals do not experience the same benefits, and every married individual may not experience the same extent of benefits. Thus, there may be wide differences in how single and married individuals behave. In particular, the context that we explore is how marriage shapes job search decisions.

%Access to some kind of insurance mechanism shapes workers' willingness to bear labor market risk \citep{landaispascalsaez,LEBARBANCHON2024435}. When workers can smooth consumption during unemployment through savings, social insurance programs, or family transfers, they can afford to remain unemployed for longer in search of a better job \citep{10.1093/qje/qjaa037}. Then, we should expect workers with greater financial security to display higher reservation wages, longer search durations, and greater mobility across jobs and sectors, possibly leading to higher quality matches. Insurance therefore plays an important role in understanding why labor market decisions get made and the timing of the decisions. While insurance is most often thought of as external to the household, marriage itself represents one of the most influential insurance arrangements in the economy. Spouses share income, consumption, assets and public goods, partially insulating one another from idiosyncratic shocks such as unemployment or illness. 

%A central insight of \cite{GULER2012352} is that couples behave differently from singles in the labor market. In their framework, married households face joint location and employment decisions that generate both new opportunities and new frictions. In particular, reservation wages for employed-searcher couples can exceed those of singles because the spouse’s income provides a form of insurance. However, while the model highlights the role of spousal income as insurance, it treats marriage as a fixed institutional arrangement. The degree of intra-household commitment and asset sharing is assumed rather than studied.

%This paper asks a more granular question: how strong is the insurance within marriage, and how does that strength shape job search behavior? We argue that the extent of marital insurance depends on the legal and institutional design of the marital contract, specifically, the matrimonial property regime. 
%Property regimes determine how assets are owned during marriage and divided upon divorce. Under separation of property, spouses retain ownership of individual assets; under community property, assets acquired during marriage are jointly owned and split equally. Universal community property extends joint ownership to premarital assets and inheritances, whereas limited community property restricts sharing to assets accumulated after marriage. These rules govern how much risk is truly pooled inside the household and, crucially, what happens in the event of divorce.

%We argue that the extent of this informal insurance depends on the strength of the marital contract, or how strong the commitment is within the relationship. \textcolor{red}{The strongest marital contract belongs to the couple who never expects to divorce, whereas weaker contracts belong to couples who internalize the possibility of divorce and prepare independently for such a possibility}. In this sense, the institutional design of divorce governs the degree of private insurance available to workers in a marriage. 

%Matrimonial property regimes determine how assets and savings are to be divided in the case of divorce, and thus they also determine the degree of shared ownership during marriage. Two broad regimes exist: separation of property, under which each spouse retains ownership of their own assets, and community property, where assets acquired during marriage are jointly owned and divided equally. Within the latter, \textit{universal} community property extends joint ownership to all premarital assets and inheritances, whereas \textit{limited} community property includes only assets acquired after marriage. These regimes have direct implications for savings and labor-supply behavior, as they dictate the extent to which partners can insure one another against income shocks \citep{voenaJMP}. 

%Unlike \cite{GULER2012352} which compares singles and couples, we compare couples facing different degrees of marital insurance. When the strength of the marriage contract declines—because fewer assets are jointly owned or because divorce reallocates assets differently—the amount of effective insurance inside marriage falls. In the event of divorce, each spouse transitions to single status with a post-divorce asset position determined by the property regime. Anticipation of this possibility affects savings, labor supply, and job search behavior even while married.
%This paper asks a more granular question: how does the strength of marital insurance, measured through the interaction between property regime and household asset structure, affect job search behavior and labor market risk-taking? 
%We exploit the 2018 Dutch reform that changed the default property regime from universal to limited community property, thereby reducing the degree of financial insurance between spouses. We combine this reform with variation in wealth shocks such as inheritances, gifts, and premarital asset appreciation. The reform re-defined which wealth shocks are shared and which remain individual. The data we use comes from the Dutch Household Survey (DHS), which records detailed information on household finances, marital status, and labor outcomes.


%Motivation
%-  \citep{GULER2012352} talks about marriage insurance and compares singles versus couples
%- We want to explicitly model the strength of the insurance WITHIN couples. Varying the marriage insurance leads to different observable job search behaviours.
%- Unlike \citep{GULER2012352}, who compare singles to couples. We're comparing COUPLES before and after their marriage strength changes to see if a decrease in marriage strengh (decrease in marriage insurance) would have different effects for men and women
%- In case of divorce, the two divorces spouses are going to be equal to singles + the assets are going to be the assets they get from divorce.

%From a theoretical viewpoint, there are additional forces that could influence joint-search decisions in the labor
%market beyond those studied in this paper. Some examples include complementarity/substitutability of leisure between
%spouses (Burdett and Mortensen, 1977), or consumption-sharing rules within the family that deviate from full income
%pooling, as in the collective model (Chiappori, 1992), or the option given to the couple to split and break up the marriage
%(Aiyagari et al., 2000), or fundamental asymmetries between men and women

%\section*{Literature Review} 
%maybe not needed for this proposal?
\section*{The Model}
As reduced form evidence, we will run a Difference-in-Differences design comparing labor market outcomes for married couples before and after the 2018 reform. This strategy identifies causal effects for already-married couples at the time of the reform, using variation in exposure induced by differences in premarital asset holdings. However, reduced form estimates have important limitations that motivate a structural approach. 

First, the reform operates through multiple simultaneous channels. A weaker property regime reduces asset pooling, alters precautionary savings motives, shifts bargaining power, and changes anticipated divorce payoffs. Reduced form estimates capture only the net effect and cannot separately identify which mechanism drives the observed changes in job search behavior. A structural model allows us to discipline these channels through the model's primitives and isolate the insurance mechanism. Second, job search decisions are dynamic. Reservation wages, unemployment durations, and on-the-job search intensity depend jointly on the value of remaining unemployed, the distribution of wage offers, the evolution of household assets, and the continuation value of marriage. The property regime affects all of these margins simultaneously. A structural model can trace these forward-looking responses and track how asset accumulation interacts with search behavior over time. Third, reduced form evidence captures local effects for couples married at the time of the reform, but leaves open the question of equilibrium responses. The reform may change which couples form, how they accumulate assets before marriage, and when they choose to divorce. These changes in the composition of couples rather than just the behavior of existing couples are outside the scope of a DiD design but are central to the long-run effects of property regime changes. A structural model allows us to incorporate endogenous marriage formation and divorce decisions, and to evaluate counterfactual policy changes beyond the specific reform we study.

%As reduced form evidence, we can run a Difference-in-Difference design to test how married couples differ in their job search behavior before and after the reform. However, reduced form won't tell us about equilibrium effects. The change in insurance may change the couple formation aspect, where couples who are formed and divorced might change due to the reform. Therefore, if we want to model these equilibrium outcomes, we need a structural framework. 

%Marital insurance through property regimes alter asset division in the event of divorce, but their effect on job search operates through forward-looking expectations about income risk, asset accumulation, and potential divorce risk. A reduced-form comparison of outcomes before and after the reform can't separately identify whether observed changes reflect shifts in wealth, divorce risk, bargaining power, or precautionary savings. A structural model allows us to discipline these channels and isolate the insurance mechanism.

%Second, job search decisions are dynamic and equilibrium objects. Reservation wages, unemployment durations, and labor market risk-taking depend on the value of remaining unemployed, the distribution of wage offers, and the continuation value of marriage. Changes in the marital contract affect all of these margins simultaneously. A structural model allows to follow the evolution of asset accumulation and run counterfactual policy analysis.








%Why do we need a structural model?

% Without a structural framework, we cannot distinguish whether weaker property regimes reduce reservation wages because insurance falls, because bargaining power shifts, or because divorce risk increases. The insurance channel operates through multiple pathways: (1) direct income pooling affects reservation wages, (2) asset sharing changes precautionary savings motives, (3) anticipated divorce outcomes alter risk-taking, and (4) spousal coordination creates strategic complementarities in job search. Reduced-form estimates give you a net effect but cannot decompose which mechanism dominates. The structural model allows you to shut down specific channels (e.g., fix divorce probabilities, eliminate savings responses, remove coordination) and quantify each mechanism's contribution

% The most compelling reason is that you want to evaluate counterfactual scenarios that don't exist in your data. 

% A reduced-form approach captures only the average effect at the time of the reform, but the structural model reveals how couples with different horizons, divorce risks, and wealth levels respond differently to the same insurance change.


%- Have two sides of the market (men and women) their decisions affect each pther
%-we have dynamics so we need to track the evolution of assets and for this we need structural model.


\subsection*{Environment}
Time is continuous and infinite. The household is a unitary decision-maker consisting of two spouses $i\in\{m,f\}$ who maximize expected discounted lifetime utility \begin{equation*}
    V_0=E_0\left[\int_0^\infty e^{\rho t}u(c_t)dt\right]
\end{equation*} where $\rho$ is the common subjective discount rate, $u(c_t)$ is the instantaneous utility from household consumption (with $u'(c)>0,\ u''(c)<0$), and $c_t$ is the household public consumption (assuming perfect pooling within marriage). 

Each spouse $i$ can be in one of two employment states, employed $(l_i=1)$ or unemployed $(l_i=0)$. The employed earn a wage $w_i$ and face an exogenous job destruction rate $\delta_i$. They participate in on-the-job search where they receive offers at rate $\gamma_i<\lambda_i$ and incur a search cost of $\kappa(s_i)$. The functional form we assume for the search cost follows from \cite{10.1093/restud/rds042} and \cite{WANG2019165}, \begin{align*}
    \kappa(s_i)=\kappa_0\cdot s_i^{\kappa_1},\quad \kappa_1>1, \quad \kappa_0 >0
\end{align*} where $s_i$ is the search intensity. This functional form captures the diminishing returns to search effort. The unemployed receive an unemployment benefit $b_i$ and search for jobs with a job offer arrival (Poisson) rate $\lambda_i$, drawn from the distribution $F_i(w)$ with support $[\underline{w}_i,\overline{w}_i]$. 

The household budget constraint takes the form \begin{equation*}
    c_t+a_{t+1}=(1+r)a_t+y_t
\end{equation*} where the household income is \begin{equation*}
    y_t=w_m\cdot\mathbf{1}\{l_m=1\}+w_f\cdot\mathbf{1}\{l_f=1\}+b_m\cdot\mathbf{1}\{l_m=0\}+b_f\cdot\mathbf{1}\{l_f=0\}.
\end{equation*}
Couples can save in a joint account $a^J$, and/or in their individual accounts $a_i$. Both type of accounts have the same risk-free rate of return $r$. The total household assets are therefore $a=a^J+a_f+a_m$. However, divorce arrives at an exogenous rate $\pi$, and the couple is subject to one of two matrimonial property regimes—universal community property or limited community property, $\theta\in\{U,L\}$. Under universal community property, it doesn't matter which account they save in, since all assets, pre-marital or otherwise, are divided equally upon divorce. In this case, we assume the couple saves in $a^J$. Under limited property, all assets accumulated prior to marriage remain separate property, and all those accumulated after marriage are joint; however, inheritance and gifts are considered separate property even if received during marriage, so long as they remain in an individual account and are not commingled (mixed with joint funds, used for joint expenses, or placed in a shared account). Let wealth shocks arrive at exogenous probability $\mu_i$, where wealth $z$ is drawn from a distribution $G_i(z)$. Universal community property was the default regime in the Netherlands until 2018, at which time the default changed to limited community property. 


\subsection*{Value Functions}
The household's state is described by $(w_m, w_f, a^J, a_m, a_f, l_m, l_f, \theta)$ where $w_i$ are wages, $a^J$ are the joint household assets, $a_i$ are individual assets, $l_i$ are employment, and $\theta$ is the marital property regime. Largely following \cite{GULER2012352} and \cite{pilossoph_wee}, there are three cases to consider:


\subsubsection*{Case 1. Dual-Unemployed, $l_m = l_f = 0$}
The household's state is described by $(a^J, a_m, a_f,\theta)$ where $a^J$ are the joint household assets, $a_i$ are individual assets, $l_i$ are employment, and $\theta$ is the marital property regime.  The value function $U(b_m,b_f,a^J, a_m, a_f,\theta)$ satisfies
\begin{align*}
    \rho U(\cdot)= \max_c\bigg\{&u(c)+\lambda_m\int_{\underline{w}}^{\overline{w}}\max\big\{W_m(w,b_f,a^J,a_m,a_f,\theta)-U(\cdot),0\big\}dF_m(w)\\&+\lambda_f\int_{\underline{w}}^{\overline{w}}\max\big\{W_f(b_m,w,a^J,a_m,a_f,\theta)-U(\cdot),0\big\}dF_f(w)\\&+\mu_m\int\big[U(b_m,b_f,a^J,a_m+z,a_f,\theta)-U(\cdot)\big]dG_m(z)\\&+\mu_f\int\big[U(b_m,b_f,a^J,a_m,a_f+z,\theta)-U(\cdot)\big]dG_f(z)\\&+\pi\left[V_m^D(\cdot)+V_f^D(\cdot)-U(\cdot)\right] + [ra + y - c]U_a(\cdot) \bigg\}.
\end{align*}

% the value of the wealth shock can be distributed according to G()
Wealth shocks arrive at rate $\mu_i$, and divorce arrives at rate $\pi$. $y$ will be equal to $b_m+ b_f$. Reservation wages $w_i^R$ are such that $W_i(w_i^R,b_i,a^J,a_m,a_f,\theta)=U(b_m,b_f,a^J,a_m,a_f,\theta)$ and are determined in equilibrium. 

\subsubsection*{Case 2. Worker-Searcher, WLOG $l_m=1,l_f=0$} 
Value function $W_m(w_m,b_f, a^J, a_m, a_f, \theta)$ satisfies \begin{align*}
    \rho W_m(\cdot) =\max_{c,s_m}\bigg\{&u(c)-\kappa(s_m)+\delta_m[U(b_m,b_f,a^J,a_m,a_f,\theta)-W_m(\cdot)]\\&+s_m\gamma_m\int_{w_m}^{\overline{w}}\big[W_m(w,b_f,a^J,a_m,a_f,\theta)-W_m(\cdot)\big]dF_m(w)\\&+\lambda_f\int_{\underline{w}}^{\overline{w}}\max\big\{E(w_m,w,a^J,a_m,a_f,\theta)-W_m(\cdot),W_f(b_m,w,a^J,a_m,a_f,\theta)-W_m(\cdot),0\big\}dF_f(w)\\&+\mu_m\int\big[W_m(w_m,b_f,a^J,a_m+z,a_f,\theta)-W_m(\cdot)\big]dG_m(z)\\&+\mu_f\int\big[W_m(w_m,b_f,a^J,a_m,a_f+z,\theta)-W_m(\cdot)\big]dG_f(z)\\&+\pi\left(V_m^D+V_f^D-W_m(\cdot)\right) + [ra + y - c] W_{a,m}(\cdot) \bigg\}.
\end{align*} The employed spouse chooses on-the-job search intensity $s_m$ to trade off between paying the search cost $\kappa(s_m)$ and having a greater probability of receiving a job offer. Upon paying the search cost $\kappa(s_m)$, they sample job offers at rate $s_m \gamma_m$. When the employed spouse gets a job offer, they accept if the offered wage is above their current wage. When the unemployed spouse gets a job offer, they can accept such that both spouses work, they can accept and the employed spouse may quit their job (the breadwinner cycle from \cite{GULER2012352}), or they can reject and remain unemployed. 

From the FOC, the optimal on-the-job search intensity is given by \begin{equation*}
    \kappa'(s_m(w_m))=\gamma_m\int_{w_m}^{\overline{w}}\big[W_m(w,b_f,a^J,a_m,a_f,\theta)-W_m(\cdot)\big]dF_m(w)
\end{equation*}
and the optimal consumption satisfies,
\[
u'(c) = W_{a,m}(\cdot)
\]

\subsubsection*{Case 3. Dual-Employed, $l_m=l_f=1$}
Value function $E(w_m,w_f,a^J,a_m,a_f,\theta)$ satisfies \begin{align*}
    \rho E(\cdot)=\max_{c,s_m,s_f}\bigg\{&u(c)-\kappa(s_m)-\kappa(s_f)\\&+\delta_m\big[W_f(b_m,w_f,a^J,a_m,a_f,\theta)-E(\cdot)\big]\\&+\delta_f\big[W_m(w_m,b_f,a^J,a_m,a_f,\theta)-E(\cdot)\big]\\&+s_m\gamma_m\int_{w_m}^{\overline{w}}\max\big\{E(w,w_f,a^J,a_m,a_f,\theta)-E(\cdot),W_m(w,b_f,a^J,a_m,a_f,\theta)-E(\cdot),0\big\}dF_m(w)\\&+s_f\gamma_f\int_{w_f}^{\overline{w}}\max\big\{E(w_m,w,a^J,a_m,a_f,\theta)-E(\cdot),W_f(b_m,w,a^J,a_m,a_f,\theta)-E(\cdot),0\big\}dF_f(w)\\&+\mu_m\int\big[E(w_m,w_f,a^J,a_m+z,a_f,\theta)-E(\cdot)\big]dG_m(z)\\&+\mu_f\int\big[E(w_m,w_f,a^J,a_m,a_f+z,\theta)-E(\cdot)\big]dG_f(z)\\&+\pi\big[V_m^D(\cdot)+V_f^D(\cdot)-E(\cdot)] + [ra + y - c] E_a(\cdot) \bigg\}
\end{align*}
\subsubsection*{Divorce Value Functions}
Upon divorce with regime $\theta$, each spouse becomes single and receives assets $a_i^D$. Under universal community property $\theta=O$, \begin{equation*}
    a_i^D=\frac{1}{2}(1-\tau_O)(a^J+a_m+a_f),
\end{equation*} and under limited community property $\theta=L$, \begin{equation*}
    a_i^D=\frac{1}{2}a^J+a_i.
\end{equation*} $\tau_\theta$ is the dissolution cost, or the fraction of assets destroyed in the divorce, and $\tau_O>\tau_L=0$. 

\paragraph{Singles' Problem}
The single-person value function for unemployment is given by $U_i(a)$, \begin{align*}
    \rho U_i^\text{sin}(a)=&\max_{c_i}\bigg\{u(c_i)+\lambda_i\int_{\underline{w}}^{\overline{w}}\max\{E_i^\text{sin}(w,a)-U_i^\text{sin}(\cdot),0\}dF_i(w) \\ &+ [ra + b_i - c_i] U_{a}^\text{sin}(a) + \eta [\chi (U(\cdot) - U_i^\text{sin})) + (1-\chi)(W_i(\cdot) - U_i^\text{sin})] \bigg\}
\end{align*}
where $[ra + b_i - c_i] U_{a}^\text{sin}(a)$ determines the additional value of the return on assets. Marriage arrives at rate $\eta$, which we assume to be exogenous and equivalent between singles and divorced individuals (for now). We assume that when $\eta$ shock hits, the couple starts at dual unemployed. The first-order condition is,
\[
u'(c_i) = U_a^\text{sin}(a)
\]

The value of employment is given by $E_i(w,a)$,
\begin{align*}
    \rho E_i^\text{sin}(\cdot)&=\max_{c_i, s_i} \bigg\{ u(c_i)-\kappa(s_i)+\delta_i[U_i^\text{sin}(a)-E_i^\text{sin}(\cdot)]\\&\quad+s_i\gamma_i\int_{w}^{\overline{w}}\big(E_i^\text{sin}(w',a)-E_i^\text{sin}(\cdot)\big)dF_i(w')\\& + [ra + w - c]E_a^\text{sin}(w,a) + \eta [\chi(E(\cdot) - E_i^\text{sin}) + (1-\chi)(W_i(\cdot) - E_i^\text{sin})] \bigg\}.
\end{align*}
We assume that once shock $\eta$ hits, the couple starts at dual employed stage. The singles' reservation wage $w_i^{R_\text{sin}}$ is such that $U_i^\text{sin}(a)=E_i^\text{sin}(w_i^{R_\text{sin}},a)$. The optimal search effort $s^*(w,a)$ is implicitly defined by \begin{align*}
    \kappa'(s_i^*(w,a))=\gamma_i\int_{w}^{\overline{w}}\big(E_i^\text{sin}(w',a)-E_i^\text{sin}(w,a)\big)dF_i(w').
\end{align*}
The optimal consumption satisfies,
\[
u'(c_i) = E_a^\text{sin}(w,a)
\]

% dont think we need the additional dissolution cost

%\section*{Model Predictions}
%\textcolor{red}{Should we try to solve the model?}
%\textcolor{red}{When insurance weakens, it depends on your asset level. If you have a lot of assets, maybe your reservation wage goes up. So maybe need to stratity by asset/wealth}

\section*{Data}
The DNB Household Survey (DHS) is a Dutch panel data set that studies the determinants of saving behavior of households, surveying more than 2,000 households each year. This dataset started in 1993, and collects information on work, pensions, housing, mortgages, income, loans, health, economic and personal characteristics.
For this project, this data has two critical advantages. First, it records granular job search behaviors of households, including self-reported reservation wages, unemployment duration, job-offer acceptance and rejection, job satisfaction, and transitions into self-employment. These measures allow us to directly examine the labor market behaviors affected by marital insurance. Second, the DHS contains detailed wealth data, including assets, liabilities, and account ownership. Importantly, savings and checking accounts specify whether they are held individually or jointly, enabling us to proxy the degree of financial pooling within households. The survey also includes information on marital status, number of children, homeownership, mortgage structure, and whether the couple is married under the default property regime or an alternative marital agreement. Some preliminary descriptive patterns from the data show that marriages under the default regime are most common, where only a small percentage of couples in our sample decide to sign a marriage settlement. This is consistent with national statistics, and helps our identification argument that marriages didn't sort into different contracts based on the reform. 

We plan to link the DHS to CBS microdata to increase statistical power of our analysis. Our identification strategy requires restricting the sample to households in which both spouses participate in the labor force and further classifying couples into dual-employed, worker-searcher, and dual-unemployed states. These restrictions substantially reduce the sample size, so the CBS data will allow us to overcome this sample issue. In addition, CBS microdata provide regional housing price indices, sector-level unemployment rates, and other local economic indicators. These variables enable us to incorporate external sources of variation in wealth shocks and labor market risk, strengthening identification of the model. Some of the datasets we can merge with CBS micro data include: 
\begin{itemize}
    \item Employment dataset: covers wages, job characteristics, job contract characteristics, unemployment insurance, dismissals
    \item Housing variation: property ownership, housing survey, housing transitions over the life cycle
    \item Household level data: household income, partner relationships, person-level register data,  household wealth (assets/debts), marriage history, partnership events
\end{itemize}

\section*{Next Steps}

Our first-pass reduced form evidence uses a Difference-in-Differences design, comparing labor market outcomes for married couples before and after the 2018 reform, exploiting variation in exposure driven by differences in premarital asset holdings. Preliminary estimates are promising but noisy, largely due to sample size constraints. Linking to CBS microdata is therefore a priority, as it will increase statistical power. Even with a larger sample, reduced form estimates capture only local average effects for couples already married at the time of the reform, and can't capture equilibrium responses. The reform may change which couples form, how individuals accumulate assets prior to marriage, and the timing of divorce, all of which affect the composition of couples observed post-reform. Addressing these equilibrium margins requires a structural framework. On the modeling side, two extensions are priorities. First, the current model takes divorce and marriage arrival rates as exogenous. In practice, both decisions are likely to depend on asset levels: individuals with greater premarital wealth may marry at different rates, and the reform changes the financial stakes of both marriage and divorce. We plan to endogenize these margins so that matching and dissolution respond to the property regime. Second, we plan to allow search intensity to vary by gender, property regime, and the asset distribution within the couple. A spouse with relatively few individual assets under limited community property faces weaker insurance and may adjust search behavior differently than their partner. 

%\textcolor{red}{We can see pre-marital assets, so we can run DiD (as a reduced form evidence). But DiD won't show up equilibrium effects (the couples who are formed, matching, divorcing might change due to the reform), they would only show local effects for already married couples}
%\textcolor{red}{Think about dndogenizing labor search intensity by gender, regime, and asset level in the couple. Lower asset person in the couple might not search as hard if they have higher insurance.}

%We're currently assuming that couples face an exogenous divorce risk, and single individuals face an exogenous marriage arrival rate. In our context, singles and couples can manipulate the rate of marriage and divorce arrival rates and we would expect these to depend on their asset levels. Singles will accumulate assets differently and get married at a different rate depending on their level of premarital wealth. Therefore, we need to model the decision to marry as an endogenous decision.

%There may be sorting of couples, different couples that may form, after regime change (depending on their pre-marital asset levels). So maybe we need to model the decision to marry.

%this would give us more power for our results since we are constrained in selecting households where both spouses are in the labor force, then selecting those in the three different categories (dual employed, worker-searcher, dual unemployed), this leaves us with small sample sizes. We can also include from the CBS micro data regional housing-price indices, sector-level unemployment rates, and local economic conditions, to incorporate external sources of variation in wealth shocks and labor-market risk.



%\begin{figure}[H]
%    \centering
%    \includegraphics[width=0.8\linewidth]{images/marital_status_composition.png}
%    \caption{This is the marital status composition. We can see that the default is the most common option in marriages.}
%\end{figure}


%\begin{figure}[H]
%    \centering
%    \includegraphics[width=0.8\linewidth]{images/married_couples_distribution.png}
%    \caption{Distribution of marriage types among married couples}
%\end{figure}



\newpage
\bibliographystyle{johd}
\bibliography{bib}



\end{document}

Event-study / staggered DiD around 2018
Inheritance event study (CBS notarial data), pre/post comparison
Financial pooling as a continuous treatment (joint-account share)
Cohabiting couples as a "pure separate" control group
Housing wealth channel via regional price appreciation
Gender asymmetry (wives vs. husbands differential response)
Sector/regional unemployment rate interaction (tests insurance demand)
Precautionary savings response (individual account balances)


DONE
% theres another term which is how your instantaenous choice of assets change the value
- another term which is how much you save
- have to take into account how your saving decision affects your value function
- derivative of the value with respect to asset times how much you're saving

DONE
% change theta not to be = L,U

DONE
% add marriage probabilities = make the simple assumption, that iin divorce they have the same probabilities. 
% if we have marriage, we can solve for steady state

% solve for the steady state
- see what happens with aggregate savings
- see comparative statics
- what happens with private savings